\section{万有形式形变与 Rim--Schlessinger 判别法}\label{src:C.4}
\begin{epigraph}
如果你无法证明自己的定理, 就不断把结论的一部分移到假设中, 直到能够证明为止.
\epigraphauthor{Ennio De Giorgi}
\end{epigraph}
\subsection{Artin 环上的函子与预叠}\label{src:C.4.1-subsection}
为简单起见, 我们在固定域 $\Bbbk$ 上工作, 但以下定义及 Rim--Schlessinger 判别法可在更一般的背景下表述 (见注记 \bookref{C.4.5}). 用 $\mathrm{Art}_{\Bbbk}$ 表示剩余域为 $\Bbbk$ 的 Artin 局部 $\Bbbk$-代数范畴.
\begin{definition}[原可表性]\label{src:C.4.1}
若存在完备 Noether 局部 $\Bbbk$-代数 $R$ 及同构 $F\xrightarrow{\sim}h_R$, 其中 $h_R:=\operatorname{Hom}_{\Bbbk\text{-}\mathrm{alg}}(R,-)$, 则称协变函子 $F:\mathrm{Art}_{\Bbbk}\to\Sets$ 为\newterm{原可表的}{prorepresentable}.
\end{definition}
若 $F:\Sch/\Bbbk\to\Sets$ 为反变函子, $x_0\in F(\Bbbk)$, 则可考虑诱导的 Artin 环函子
\[
F_{x_0}:\mathrm{Art}_{\Bbbk}\to\Sets,\qquad A\mapsto\{x\in F(A)\mid x|_{\Bbbk}=x_0\in F(\Bbbk)\},
\]
其中 $x|_{\Bbbk}$ 表示 $x$ 在 $F(A)\to F(A/\mathfrak m_A)$ 下的像. 若 $F$ 可由概形 $X$ 表示, $x\in X$ 为对应 $x_0$ 的 $\Bbbk$-点, 则 $F_{x_0}$ 可由 $\widehat{\mathcal O}_{X,x}$ 原表示. 可能 $F_{x_0}$ 原可表, 但 $F$ 不可表.

许多 Artin 环函子并非原可表. 例如, 若 $C_0$ 为具有非平凡自同构群的光滑连通射影曲线, 则协变函子 $F_{C_0}:\mathrm{Art}_{\Bbbk}\to\Sets$ 不是原可表的; 其中 $F_{C_0}(A)$ 由满足 $\mathcal C\times_AA/\mathfrak m_A$ 与 $C_0$ 同构的光滑固有曲线族 $\mathcal C\to\Spec A$ 的同构类组成. 然而, 许多模函子具有\textit{万有}形变. 由于记录自同构很重要, 我们将对 $\mathrm{Art}_{\Bbbk}^{\mathrm{op}}$ 上的预叠表述定义.
\begin{definition}[万有性]\label{src:C.4.2}
设 $\mathcal X$ 为 $\mathrm{Art}_{\Bbbk}^{\mathrm{op}}$ 上的预叠, 使得群胚 $\mathcal X(\Bbbk)$ 等价于集合 $\{x_0\}$.
\begin{enumerate}[(1)]
\item $x_0$ 的\newterm{形式形变}{formal deformation} $(R,\{x_n\})$ 是完备 Noether 局部 $\Bbbk$-代数 $(R,\mathfrak m_R)$, 连同对象 $x_n\in\mathcal X(R/\mathfrak m_R^{n+1})$ 及 $\Spec R/\mathfrak m_R^n\to\Spec R/\mathfrak m_R^{n+1}$ 上的态射 $x_{n-1}\to x_n$ 的资料.
\item 若对 $\mathrm{Art}_{\Bbbk}$ 中的每个满射 $A\twoheadrightarrow A_0$, 满足 $\mathfrak m_A^{n+1}=0$, 每个对象 $\eta\in\mathcal X(A)$, 以及 $\Bbbk$-代数同态 $\phi_0:R/\mathfrak m_R^{n+1}\to A_0$ 和 $\mathcal X(A_0)$ 中的同构 $\alpha_0:x_n|_{A_0}\xrightarrow{\sim}\eta|_{A_0}$, 都存在延拓 $\phi_0$ 的 $\Bbbk$-代数同态 $\phi:R/\mathfrak m_R^{n+1}\to A$ 及延拓 $\alpha_0$ 的 $\mathcal X(A)$ 中的同构 $\alpha:x_n|_A\xrightarrow{\sim}\eta$, 则称形式形变 $(R,\{x_n\})$ 为\newterm{万有的}{versal}.
\item 若万有形式形变 $(R,\{x_n\})$ 在同构类上诱导的映射 $h_R(\Bbbk[\epsilon])\to\mathcal X(\Bbbk[\epsilon])/{\sim}$ 双射, 则称其为\newterm{极小万有}{miniversal}的 (或\newterm{原可表包络}{prorepresentable hull}).
\end{enumerate}
\end{definition}
换言之, 形式形变是 $\lim\mathcal X(R/\mathfrak m^n)$ 的元素. 当 $\mathcal X=F$ 为协变函子 $\mathrm{Art}_{\Bbbk}\to\Sets$ 时, 形式形变是一列相容元素 $x_n\in F(R/\mathfrak m_R^{n+1})$. 若 $F$ 可由 $R$ 原表示, $x_n\in F(R/\mathfrak m_R^{n+1})$ 为对应元素, 则 $\{x_n\}$ 为极小万有形式形变, 此时 \bookeqref{C.4.3} 中存在唯一提升. 如同形变 $x_n\in\mathcal X(R/\mathfrak m_R^{n+1})$ 可由 Yoneda 的 $2$-引理看成态射 $\Spec R/\mathfrak m_R^{n+1}\to\mathcal X$, 形式形变可看成预叠态射 $\{x_n\}:h_R\to\mathcal X$. 从这一角度看, $\{x_n\}$ 万有, 当且仅当对每个交换图
\begin{equation}\label{src:C.4.3}
\begin{tikzcd}
\Spec A_0\arrow[r]\arrow[d,hook]&h_R\arrow[d,"{\{x_n\}}"]\\
\Spec A\arrow[r,"\eta"']\arrow[ur,dashed]&\mathcal X,
\end{tikzcd}
\end{equation}
都存在提升, 其中 $A\twoheadrightarrow A_0$ 为 $\mathrm{Art}_{\Bbbk}$ 中的满射. 应将其与光滑性的无穷小提升判据 (\bookref{A.3.1} 和 \bookref{4.7.1}) 比较. 注意每个满射 $A\twoheadrightarrow A_0$ 都分解为一维核的满射的复合, 所以可在 $\ker(A\to A_0)\simeq\Bbbk$ 的满射上检验万有性. 万有形变极小万有, 当且仅当它在切空间上诱导同构 $h_R(\Bbbk[\epsilon])\to\mathcal X(\Bbbk[\epsilon])/{\sim}$.
\begin{remark}[从整体预叠到局部形变预叠]\label{src:C.4.4}
若 $\mathcal X$ 为 $\Sch/\Bbbk$ 上的预叠, $x_0\in\mathcal X(\Bbbk)$, 则定义 $x_0$ 处的\newterm{局部形变预叠}{local deformation prestack} $\mathcal X_{x_0}$ 如下: 一个对象是 $\mathcal X$ 中的态射 $\alpha:x_0\to x$, 其中 $x$ 为环 $A\in\mathrm{Art}_{\Bbbk}^{\mathrm{op}}$ 上的对象; 换言之, 对象为二元组 $(x,\alpha)$, 其中 $x\in\mathcal X(A)$, $\alpha:x_0\to x|_{\Bbbk}$ 为同构. 态射 $(\alpha:x_0\to x)\to(\alpha':x_0\to x')$ 是满足 $\alpha'=\alpha\circ\beta$ 的态射 $\beta:x\to x'$. 注意纤维范畴 $\mathcal X_{x_0}(\Bbbk)$ 等价于集合 $\{\mathrm{id}:x_0\to x_0\}$.

若 $\mathcal X$ 为代数叠, 有来自 $\Bbbk$-概形 $U$ 的光滑呈示 $U\to\mathcal X$, 且 $u\in U(\Bbbk)$ 为 $x_0$ 的提升, 则将 $x_n\in\mathcal X(\mathcal O_{U,u}/\mathfrak m_u^{n+1})$ 取为复合 $\Spec\mathcal O_{U,u}/\mathfrak m_u^{n+1}\to U\to\mathcal X$, 就定义万有形式形变 $\{x_n\}$. Rim--Schlessinger 判别法 (定理 \bookref{C.4.6}) 给出万有形式形变存在的判据, Artin 代数性公理 (定理 \bookref{C.7.1}) 则进一步给出万有形式形变由上述光滑呈示 $U\to\mathcal X$ 诱导的判据.
\end{remark}
\begin{remark}\label{src:C.4.5}
为在更一般的基底上工作, 考虑剩余域为 $\Bbbk$ 的完备 Noether 局部环 $\Lambda$, 并将 $\mathrm{Art}_\Lambda$ 定义为带有等同 $\Bbbk\xrightarrow{\sim}A/\mathfrak m$ 的 Artin 局部 $\Lambda$-代数 $(A,\mathfrak m)$ 范畴. 实际上, 常取 $\Lambda$ 为 Witt 向量环, 例如 $\Lambda=\mathbb Z_p$. 这种一般性对许多应用很重要, 例如将对象从特征 $p$ 提升到特征 $0$; 见 \S\bookref{C.5.3-subsection}. 更一般地, 还可考虑 $A\to\Bbbk$ 有限 (但未必满射) 的背景; 见 \cite[Tag 06GB]{SP}.
\end{remark}
\subsection{Rim--Schlessinger 判别法}\label{src:C.4.2-subsection}
Rim--Schlessinger 判别法给出 $\mathrm{Art}_{\Bbbk}^{\mathrm{op}}$ 上的预叠 $\mathcal X$ (或 Schlessinger 原始表述中的协变函子 $F:\mathrm{Art}_{\Bbbk}\to\Sets$) 具有万有形式形变的充要条件.
\begin{theorem}[Rim--Schlessinger 判别法]\label{src:C.4.6}
设 $\mathcal X$ 为 $\mathrm{Art}_{\Bbbk}^{\mathrm{op}}$ 上的预叠, 群胚 $\mathcal X(\Bbbk)$ 等价于集合 $\{x_0\}$. 对 $\mathrm{Art}_{\Bbbk}$ 中的态射 $B_0\to A_0$ 和 $A\to A_0$, 考虑自然函子
\begin{equation}\label{src:C.4.7}
\mathcal X(B_0\times_{A_0}A)\to\mathcal X(B_0)\times_{\mathcal X(A_0)}\mathcal X(A).
\end{equation}
则 $\mathcal X$ 具有极小万有形式形变, 当且仅当
\begin{enumerate}[$(\mathrm{RS}_1)$]
\item 当 $A\twoheadrightarrow A_0$ 为核 $\Bbbk$ 的满射时, 函子 \bookeqref{C.4.7} 本质满射;
\item 当 $A_0=\Bbbk$ 且 $A=\Bbbk[\epsilon]$ 时, 映射 \bookeqref{C.4.7} 本质满射, 且给定两个交换图
\[
\begin{tikzcd}[column sep=small]
x_0\arrow[r]\arrow[d]&y_0\arrow[d,"\alpha"]\arrow[dr,"\alpha'"]&\\
x\arrow[r,"\gamma"]\arrow[rr,bend right=25,"\gamma'"']&y\arrow[r,dashed,"\beta"]&y'
\end{tikzcd}
\qquad\text{在}\qquad
\begin{tikzcd}[column sep=small]
\Spec\Bbbk\arrow[r,hook]\arrow[d]&\Spec B_0\arrow[d]\\
\Spec\Bbbk[\epsilon]\arrow[r,hook]&\Spec(\Bbbk[\epsilon]\times_{\Bbbk}B_0),
\end{tikzcd}
\]
都存在 $\mathcal X(\Bbbk[\epsilon]\times_{\Bbbk}B_0)$ 中的同构 $\beta:y\to y'$, 满足 $\alpha'=\beta\circ\alpha$ (但不要求 $\gamma'=\gamma\circ\beta$).
\item $\dim_{\Bbbk}T_{\mathcal X}<\infty$, 其中 $T_{\mathcal X}:=\mathcal X(\Bbbk[\epsilon])/{\sim}$ (由引理 \bookref{C.4.9}, 它为向量空间).
\end{enumerate}
此外, $\mathcal X$ 原可表, 当且仅当 $\mathcal X$ 等价于函子, 且
\begin{enumerate}[$(\mathrm{RS}_1)$]\setcounter{enumi}{3}
\item 当 $A\twoheadrightarrow A_0$ 为核 $\Bbbk$ 的满射时, 映射 \bookeqref{C.4.7} 为等价.
\end{enumerate}
\end{theorem}
条件 $(\mathrm{RS}_2)$--$(\mathrm{RS}_3)$ (有时称为\newterm{半齐性}{semi-homogeneity}) 可能不易理解, 但实际中验证更强的条件 $(\mathrm{RS}_4)$ (称为\newterm{齐性}{homogeneity}) 常常同样容易, 甚至验证更强的条件 $(\mathrm{RS}_4^*)$ (称为\newterm{强齐性}{strong homogeneity}) 也是如此; 见 \S\bookref{C.4.3-subsection}. 当 $A\to A_0$ 满射时, $\Spec(B_0\times_{A_0}A)$ 为概形范畴中 $\Spec A_0\hookrightarrow\Spec A$ 与 $\Spec B_0\to\Spec B$ 的推出 (见 \S\bookref{B.4}). 因此, 齐性条件变为推出上对象的粘合条件.
\begin{remark}[Schlessinger 判别法]\label{src:C.4.8}
当 $\mathcal X$ 为协变函子 $F:\mathrm{Art}_{\Bbbk}\to\Sets$, 且 $F(\Bbbk)=\{x_0\}$ 时, $(\mathrm{RS}_1)$--$(\mathrm{RS}_4)$ 化为 \cite{Sch68} 所引入的 Schlessinger 条件:
\begin{enumerate}[$(\mathrm H_1)$]
\item 当 $A\twoheadrightarrow A_0$ 为核 $\Bbbk$ 的满射时, 映射 \bookeqref{C.4.7} 满射;
\item 当 $A_0=\Bbbk$ 且 $A=\Bbbk[\epsilon]$ 时, 映射 \bookeqref{C.4.7} 双射;
\item $\dim_{\Bbbk}F(\Bbbk[\epsilon])<\infty$; 且
\item 当 $A\twoheadrightarrow A_0$ 为核 $\Bbbk$ 的满射时, 映射 \bookeqref{C.4.7} 双射.
\end{enumerate}
若 $(\mathrm H_1)$--$(\mathrm H_3)$ 成立, 则函子 $F$ 具有极小万有形式形变; 若 $(\mathrm H_3)$--$(\mathrm H_4)$ 成立, 则 $F$ 原可表.

若 $\mathrm{Art}_{\Bbbk}$ 上的预叠 $\mathcal X$ 满足 $(\mathrm{RS}_1)$--$(\mathrm{RS}_3)$, 则参数化对象同构类的函子 $F_{\mathcal X}:\mathrm{Art}_{\Bbbk}\to\Sets$ 满足 $(\mathrm H_1)$--$(\mathrm H_3)$, 但反之未必成立. 此外, $\mathcal X(B_0\times_{A_0}A)\to\mathcal X(B_0)\times_{\mathcal X(A_0)}\mathcal X(A)$ 的本质满射性蕴含 $F_{\mathcal X}(B_0\times_{A_0}A)\to F_{\mathcal X}(B_0)\times_{F_{\mathcal X}(A_0)}F_{\mathcal X}(A)$ 的满射性; 只要对对象 $y_0\in\mathcal X(B_0)$, $\operatorname{Aut}_{\mathcal X(B_0)}(y_0)\to\operatorname{Aut}_{\mathcal X(A_0)}(y_0|_{A_0})$ 满射, 则 $\mathcal X$ 的全忠实性蕴含 $F_{\mathcal X}$ 的单射性. 后一条件在 $F_{\mathcal X}(A_0)$ 为集合时成立, 例如 $A_0=\Bbbk$ 时. 若如注记 \bookref{C.4.4}, $\mathcal X$ 为 $\Sch/\Bbbk$ 上代数叠 $\widetilde{\mathcal X}$ 的对象 $x_0\in\widetilde{\mathcal X}(\Bbbk)$ 所产生的局部形变预叠, 则由光滑性的无穷小提升判据 (\bookref{4.7.1}), 自同构上的满射性条件变为惯性叠 $I_{\mathcal X}\to\mathcal X$ 在 $e(x_0)$ 处的光滑性, 其中 $e:\mathcal X\to I_{\mathcal X}$ 为恒等截面.

虽然 $F_{\mathcal X}$ 的极小万有形式形变的存在已足以满足许多应用, 但对于带自同构的模问题, 要求 $\mathcal X$ 的极小万有形式形变存在更自然, 且某些应用确实需要这一一般性, 例如 Artin 代数化 (定理 \bookref{C.6.8}) 和 Artin 代数性公理 (定理 \bookref{C.7.4}).
\end{remark}
开始证明前, 先说明条件 $(\mathrm{RS}_1)$--$(\mathrm{RS}_4)$ 的一些性质.
